% !PRINTED 703
\weberpage{703}
%<<CONT
gives a contribution to the term with $z_1^{-1}$, and by (2) this has the coefficient $a_{-1}$. From this it follows:

2. \gesperrt{The residue of $dJ$ is independent of the choice of the function $z$.}

The following theorem holds further:

3. \gesperrt{The sum of all residues of a differential is equal to zero.}

We extend the expression $\mathfrak{A}/\mathfrak{B}$ of the differential $dJ$, if necessary, by the addition of points in the numerator and denominator, so that the mutually distinct points of the lower polygon
\weq{3}{\mathfrak{P}_1 \mathfrak{P}_2 \ldots \mathfrak{P}_n}
form a polygon of a proper class. The poles of $dJ$ are then contained among these $\mathfrak{P}$, and possibly other points as well, for which the residue of $dJ$ is equal to zero.

As denominator of $dJ$ we may then take
\weq{4}{\mathfrak{B} = \mathfrak{P}_1^{m_1} \mathfrak{P}_2^{m_2} \ldots \mathfrak{P}_n^{m_n}}
if we add the points taken in excess back again in the numerator.

We take a function $z$ in $\Omega$ of the $n$th order, which becomes infinite at each of the points $\mathfrak{P}_i$ and at each of them only to the first order, and form the development for the point $\mathfrak{P}_i$:
\weq{5}{\frac{dJ}{dz} = a_{m-2}^{(i)}z^{m-2} + \cdots + a_0^{(i)} + a_{-1}^{(i)}z^{-1} + \eta^{(i)}z^{-2}.}

If we admit the value 0 for the coefficient $a_r^{(i)}$, we can begin at all the points $\mathfrak{P}_i$ with the same power $z^{m-2}$. The sum of the residues of $dJ$ is
\[
-\overset{i}{\sum} a_{-1}^{(i)}.
\]

According to § 182, 11. a system of functions $\alpha_1, \alpha_2, \ldots, \alpha_n$ can be determined in such a way that $\alpha_r$ becomes infinitely small in the first order at the point $\mathfrak{P}_r$, while at the remaining points $\mathfrak{P}_i$ it remains finite and different from zero.

If one then puts
\[
\alpha_1^m \alpha_2^m \ldots \alpha_n^m = \alpha_i^m \varrho_i,
\]
there results
\weq{6}{\begin{aligned}
\varrho_1 &= 0^m \text{ in } \mathfrak{P}_2, \mathfrak{P}_3, \ldots, \mathfrak{P}_n, & \text{finite and not} &= 0 \text{ in } \mathfrak{P}_1, \\
\varrho_2 &= 0^m \text{ ,, } \mathfrak{P}_1, \mathfrak{P}_3, \ldots, \mathfrak{P}_n, & \text{,, ,, ,,} &= 0 \text{ ,, } \mathfrak{P}_2, \\
&\cdots\cdots\cdots\cdots\cdots\cdots\cdots\cdots \\
\varrho_n &= 0^m \text{ ,, } \mathfrak{P}_1, \mathfrak{P}_2, \ldots, \mathfrak{P}_{n-1} & \text{,, ,, ,,} &= 0 \text{ in } \mathfrak{P}_n.
\end{aligned}}
